\begin{abstract}
Geometric calibration makes thermal models useful for laser processing, but a fitted melt boundary leaves the underlying cooling history dependent on source and material assumptions. Here, a nonlinear enthalpy conduction model is used to examine the thermal interpretation of length-calibrated bare-plate IN625 tracks. Conductivity and sensible heat-capacity continuations above their tabulated ranges are varied independently while the calibrated source and phase law remain fixed. Matching length leaves a wider and shallower fusion zone than the reference means. Within the controlled comparison, conductivity holding extends the pool, whereas heat-capacity holding shortens it. The conductivity-induced extension mainly displaces the rear solidus and liquidus together, with much less change in their separation. Across two scan speeds, nearly unchanged rear spatial phase spans correspond to a material passage time shorter by about one-third at the higher speed. These results distinguish translation of the thermal tail from enlargement of the phase region and connect spatial geometry to the time experienced by material. They provide a basis for interpreting property sensitivity after geometric calibration and for selecting thermal observations that assess the resulting histories.
\end{abstract}
